%%%%%%%%%%%%%%%%%%%%%%%%%%%%%%%%%%%%%%%%%
% kaobook
% LaTeX Template
% Version 1.2 (4/1/2020)
%
% This template originates from:
% https://www.LaTeXTemplates.com
%
% For the latest template development version and to make contributions:
% https://github.com/fmarotta/kaobook
%
% Authors:
% Federico Marotta (federicomarotta@mail.com)
% Based on the doctoral thesis of Ken Arroyo Ohori (https://3d.bk.tudelft.nl/ken/en)
% and on the Tufte-LaTeX class.
% Modified for LaTeX Templates by Vel (vel@latextemplates.com)
%
% License:
% CC0 1.0 Universal (see included MANIFEST.md file)
%
%%%%%%%%%%%%%%%%%%%%%%%%%%%%%%%%%%%%%%%%%

%----------------------------------------------------------------------------------------
%	PACKAGES AND OTHER DOCUMENT CONFIGURATIONS
%----------------------------------------------------------------------------------------

\documentclass[
	fontsize=13pt, % Base font size
	twoside=false, % Use different layouts for even and odd pages (in particular, if twoside=true, the margin column will be always on the outside)
	%open=any, % If twoside=true, uncomment this to force new chapters to start on any page, not only on right (odd) pages
	%chapterprefix=true, % Uncomment to use the word "Chapter" before chapter numbers everywhere they appear
	%chapterentrydots=true, % Uncomment to output dots from the chapter name to the page number in the table of contents
	numbers=noenddot, % Comment to output dots after chapter numbers; the most common values for this option are: enddot, noenddot and auto (see the KOMAScript documentation for an in-depth explanation)
	%draft=true, % If uncommented, rulers will be added in the header and footer
	%overfullrule=true, % If uncommented, overly long lines will be marked by a black box; useful for correcting spacing problems
]{kaobook}
%!TEX root = ../thesis.tex

%\hypersetup{colorlinks,linktocpage,urlcolor=red}
%
%\definecolor{myGreen}{HTML}{05C18E} 
%\definecolor{myGreenDarker}{HTML}{178C6C} \colorlet{mylinkcolor}{green!50!black}

\definecolor{webbrown}{rgb}{.6,0,0}

\hypersetup{
  colorlinks=true,
  linkcolor=black, %myGreenDarker
%  urlcolor=myGreenDarker,
  citecolor = webbrown,
  urlcolor=webbrown,
  hyperfootnotes=false,
  hypertexnames,
  bookmarks=true}
  
%\setsidenotefont{\color{black}\footnotesize}   <-- set the color and font here
%\setmarginnotefont{\color{black}\footnotesize} <-- and here
%
%\renewcommand{\maketitlepage}[0]{%
%  \cleardoublepage%
%  {%
%  \sffamily%
%  \begin{fullwidth}%
%  \fontsize{18}{20}\selectfont\par\noindent\textcolor{darkgray}{\allcaps{\thanklessauthor}}%
%  \vspace{11.5pc}%
%  \fontsize{24}{45}\selectfont\par\noindent\textcolor{darkgray}{\allcaps{\thanklesstitle}}
%  \fontsize{17.4}{25}\selectfont\par\noindent\textcolor{darkgray}{\allcaps{For Affective Touch Communication}}%
%  \fontsize{10.0}{17}\selectfont\par\noindent\textcolor{Gray}{\allcaps{Devices that touch to convey emotions and feel that contact}}%
%
%  \vfill%
%  \fontsize{14}{16}\selectfont\par\noindent\allcaps{\thanklesspublisher}%
%  \end{fullwidth}%
%  }
 %  \thispagestyle{empty}%
%  \clearpage%
%}

%%%% Kevin Godny's code for title page and contents from https://groups.google.com/forum/#!topic/tufte-latex/ujdzrktC1BQ
% \makeatletter
% \renewcommand{\maketitlepage}{%
% \begingroup%
% \setlength{\parindent}{0pt}

% {\fontsize{24}{24}\selectfont\textit{\@author}\par}

% \vspace{1.75in}{\fontsize{36}{54}\selectfont\@title\par}

% % \vspace{0.5in}{\fontsize{14}{14}\selectfont\textsf{\smallcaps{\@date}}\par}
% \vspace{0.5in}{\fontsize{14}{14}\selectfont\textsf{\sc\@date}\par}

% \vfill{\fontsize{14}{14}\selectfont\textit{\@publisher}\par}

% \thispagestyle{empty}
% \endgroup
% }
% \makeatother

%\titlecontents{part}%
%    [0pt]% distance from left margin
%    {\addvspace{0.25\baselineskip}}% above (global formatting of entry)
%    {\allcaps{Part~\thecontentslabel}\allcaps}% before w/ label (label = ``Part I'')
%    {\allcaps{Part~\thecontentslabel}\allcaps}% before w/o label
%    {}% filler and page (leaders and page num)
%    [\vspace*{0.5\baselineskip}]% after
%
%
%\titlecontents{chapter}%
%    [4em]% distance from left margin
%    {}% above (global formatting of entry)
%    {\contentslabel{2em}\textit}% before w/ label (label = ``Chapter 1'')
%    {\hspace{0em}\textit}% before w/o label
%    {\qquad\thecontentspage}% filler and page (leaders and page num)
%    [\vspace*{0.5\baselineskip}]% after
%%%%% End additional code by Kevin Godby


%% CHANGE CITE COMMAND
\renewcommand{\cite}[1]{%
~\citep{#1}%
}



%%
% If they're installed, use Bergamo and Chantilly from www.fontsite.com.
% They're clones of Bembo and Gill Sans, respectively.
%\IfFileExists{bergamo.sty}{\usepackage[osf]{bergamo}}{}% Bembo
%\IfFileExists{chantill.sty}{\usepackage{chantill}}{}% Gill Sans

%%%%%%%%%%%%%%%%%%%%%%%%%%%%%%%%%%%%%%%%%%%%%%%%%%%%%%%%%%
%%% INCLUSION / EXCLUSION %%%%%%%%%%%%%%%%%%%
\usepackage{microtype}
\usepackage{comment}
% !!! Comment or uncomment line under to exclude or include the content of the chapter:
%\excludecomment{content} % exclude the content, (only get introduction and summary)
\includecomment{content} % include the content, (get eveevolutionrything)
\includecomment{export}
%%%%%%%%%%%%%%%%%%%%%%%%%%%%%%%%%%%%%%%%%%%%%%%%%%%%%%%%%%


%%
% For nicely typeset tabular material
\usepackage{booktabs}
%%
% For graphics / images
\usepackage{graphicx}
\setkeys{Gin}{width=\linewidth,totalheight=\textheight,keepaspectratio}
\graphicspath{{graphics/}}
% The fancyvrb package lets us customize the formatting of verbatim environments.  We use a slightly smaller font.
\usepackage{fancyvrb}
\fvset{fontsize=\normalsize}

%%
% Prints argument within hanging parentheses (i.e., parentheses that take
% up no horizontal space).  Useful in tabular environments.
% \newcommand{\hangp}[1]{\makebox[0pt][r]{(}#1\makebox[0pt][l]{)}}

%%
% Prints an asterisk that takes up no horizontal space.
% Useful in tabular environments.
% \newcommand{\hangstar}{\makebox[0pt][l]{*}}

%%
% Prints a trailing space in a smart way.
\usepackage{xspace}

%
%%%
%% Some shortcuts for Tufte's book titles.  The lowercase commands will
%% produce the initials of the book title in italics.  The all-caps commands
%% will print out the full title of the book in italics.
%\newcommand{\vdqi}{\textit{VDQI}\xspace}
%\newcommand{\ei}{\textit{EI}\xspace}
%\newcommand{\ve}{\textit{VE}\xspace}
%\newcommand{\be}{\textit{BE}\xspace}
%%\newcommand{\VDQI}{\textit{Visualizing dynamic social data  with rationally designed constructive systems}\xspace}
%\newcommand{\EI}{\textit{Envisioning Information}\xspace}
%\newcommand{\VE}{\textit{Visual Explanations}\xspace}
%\newcommand{\BE}{\textit{Beautiful Evidence}\xspace}
%\newcommand{\TL}{Tufte-\LaTeX\xspace}

% Prints the month name (e.g., January) and the year (e.g., 2008)
% \newcommand{\monthyear}{%
%   \ifcase\month\or January\or February\or March\or April\or May\or June\or
%   July\or August\or September\or October\or November\or
%   December\fi\space\number\year
% }





% Prints an epigraph and speaker in sans serif, all-caps type.
\newcommand{\openepigraph}[2]{%
  %\sffamily\fontsize{14}{16}\selectfont
  \begin{fullwidth}
  \sffamily\large
  \begin{doublespace}
  \noindent\allcaps{#1}\\% epigraph
  \noindent\allcaps{#2}% author
  \end{doublespace}
  \end{fullwidth}
}

% Inserts a blank page
% \newcommand{\blankpage}{\newpage\hbox{}\thispagestyle{empty}\newpage}

\usepackage{units}

% Typesets the font size, leading, and measure in the form of 10/12x26 pc.
\newcommand{\measure}[3]{#1/#2$\times$\unit[#3]{pc}}

% Macros for typesetting the documentation
\newcommand{\hlred}[1]{\textcolor{Green}{#1}}% prints in red
\newcommand{\hangleft}[1]{\makebox[0pt][r]{#1}}
% \newcommand{\hairsp}{\hspace{1pt}}% hair space
\newcommand{\hquad}{\hskip0.5em\relax}% half quad space
\newcommand{\TODO}{\textcolor{red}{\bf TODO!}\xspace}
% \newcommand{\ie}{\textit{i.\hairsp{}e.}\xspace}
% \newcommand{\eg}{\textit{e.\hairsp{}g.}\xspace}
% \newcommand{\na}{\quad--}% used in tables for N/A cells
\providecommand{\XeLaTeX}{X\lower.5ex\hbox{\kern-0.15em\reflectbox{E}}\kern-0.1em\LaTeX}
\newcommand{\tXeLaTeX}{\XeLaTeX\index{XeLaTeX@\protect\XeLaTeX}}
% \index{\texttt{\textbackslash xyz}@\hangleft{\texttt{\textbackslash}}\texttt{xyz}}
\newcommand{\tuftebs}{\symbol{'134}}% a backslash in tt type in OT1/T1
\newcommand{\doccmdnoindex}[2][]{\texttt{\tuftebs#2}}% command name -- adds backslash automatically (and doesn't add cmd to the index)
\newcommand{\doccmddef}[2][]{%
  \hlred{\texttt{\tuftebs#2}}\label{cmd:#2}%
  \ifthenelse{\isempty{#1}}%
    {% add the command to the index
      \index{#2 command@\protect\hangleft{\texttt{\tuftebs}}\texttt{#2}}% command name
    }%
    {% add the command and package to the index
      \index{#2 command@\protect\hangleft{\texttt{\tuftebs}}\texttt{#2} (\texttt{#1} package)}% command name
      \index{#1 package@\texttt{#1} package}\index{packages!#1@\texttt{#1}}% package name
    }%
}% command name -- adds backslash automatically
\newcommand{\doccmd}[2][]{%
  \texttt{\tuftebs#2}%
  \ifthenelse{\isempty{#1}}%
    {% add the command to the index
      \index{#2 command@\protect\hangleft{\texttt{\tuftebs}}\texttt{#2}}% command name
    }%
    {% add the command and package to the index
      \index{#2 command@\protect\hangleft{\texttt{\tuftebs}}\texttt{#2} (\texttt{#1} package)}% command name
      \index{#1 package@\texttt{#1} package}\index{packages!#1@\texttt{#1}}% package name
    }%
}% command name -- adds backslash automatically
\newcommand{\docopt}[1]{\ensuremath{\langle}\textrm{\textit{#1}}\ensuremath{\rangle}}% optional command argument
\newcommand{\docarg}[1]{\textrm{\textit{#1}}}% (required) command argument
\newenvironment{docspec}{\begin{quotation}\ttfamily\parskip0pt\parindent0pt\ignorespaces}{\end{quotation}}% command specification environment
\newcommand{\docenv}[1]{\texttt{#1}\index{#1 environment@\texttt{#1} environment}\index{environments!#1@\texttt{#1}}}% environment name
\newcommand{\docenvdef}[1]{\hlred{\texttt{#1}}\label{env:#1}\index{#1 environment@\texttt{#1} environment}\index{environments!#1@\texttt{#1}}}% environment name
\newcommand{\docpkg}[1]{\texttt{#1}\index{#1 package@\texttt{#1} package}\index{packages!#1@\texttt{#1}}}% package name
\newcommand{\doccls}[1]{\texttt{#1}}% document class name
\newcommand{\docclsopt}[1]{\texttt{#1}\index{#1 class option@\texttt{#1} class option}\index{class options!#1@\texttt{#1}}}% document class option name
\newcommand{\docclsoptdef}[1]{\hlred{\texttt{#1}}\label{clsopt:#1}\index{#1 class option@\texttt{#1} class option}\index{class options!#1@\texttt{#1}}}% document class option name defined
\newcommand{\docmsg}[2]{\bigskip\begin{fullwidth}\noindent\ttfamily#1\end{fullwidth}\medskip\par\noindent#2}
\newcommand{\docfilehook}[2]{\texttt{#1}\index{file hooks!#2}\index{#1@\texttt{#1}}}
\newcommand{\doccounter}[1]{\texttt{#1}\index{#1 counter@\texttt{#1} counter}}




%\geometry{textwidth=.55\paperwidth}


% Generates the index
\usepackage{makeidx}
\makeindex



% Nomenclature
%\usepackage{nomencl}
%\renewcommand{\nomname}{List of Abbreviations}
%\makenomenclature



%%%%
\makeatletter
\renewcommand*\l@figure{\@dottedtocline{1}{1.5em}{2.3em}}
\makeatother

%% change TOC
\setcounter{tocdepth}{2}
\setcounter{secnumdepth}{2}

%%%%%%%%%%%%%%%%%%%%%%%%%%%%%%%%%%%%%%%%%%%%%%%%%%
%%%%%%%%%%%%%%%%%%%%%%%%%%%%%%%%%%%%%%%%%%%%%%%%%%
\usepackage{amssymb}% http://ctan.org/pkg/amssymb
\usepackage{pifont}% http://ctan.org/pkg/pifont
%\usepackage{graphics} % for EPS, load graphicx instead
\usepackage{graphicx}
\usepackage{multirow}
\usepackage{xspace}
\usepackage{tabularx}
\usepackage{color}
\usepackage{listings}
\usepackage{ulem}
\usepackage{colortbl}
\usepackage{morefloats}
\usepackage{enumitem}
\usepackage{rotating}
\usepackage{comment}
\usepackage{rotating}
% \usepackage[sort, numbers]{natbib} 
\usepackage[retainorgcmds]{IEEEtrantools}
\usepackage{bibentry}
\usepackage{longtable}
\usepackage{glossaries}
\usepackage{gensymb}
\usepackage{csvsimple}
\usepackage{amsmath}
\usepackage{cleveref}% Has to be loaded after hyperref
\usepackage[utf8]{inputenc}
\usepackage{todonotes}
\usepackage{marginfix}
\usepackage[export]{adjustbox}
\usepackage{fullwidth}
%\setkeys{Gin}{height=2cm}
%\usepackage{float}
\usepackage[caption=false]{subfig}

\usepackage[strict]{changepage}

\setlist[description]{style = multiline, labelwidth = 55pt}
\usepackage[parfill]{parskip}
\makeatletter
% Paragraph indentation and separation for normal text
% \renewcommand{\@tufte@reset@par}{%
%   \setlength{\RaggedRightParindent}{1.0pc}%
%   \setlength{\parindent}{1pc}%
%   \setlength{\parskip}{8pt}%\baselineskip % default 12pt for 10pt font
% }
% \@tufte@reset@par

% Paragraph indentation and separation for marginal text
% \renewcommand{\@tufte@margin@par}{%
%   \setlength{\RaggedRightParindent}{0.5pc}%
%   \setlength{\JustifyingParindent}{0.0pc}%
%   \setlength{\parindent}{0.5pc}%
%   \setlength{\parskip}{6pt}%
% }
\makeatother





%% Correction

%\newcommand{\Ssubsection}[1]{{\setlength{\parindent}{0cm}\normalfont{\textit{\newline#1}}}\newline}
%\newcommand{\Ssubsection}[1]{{\setlength{\parindent}{0cm}\normalfont{\textit{#1}}}}




\usepackage{mdframed}
\newmdenv[
  leftmargin = 0pt,
  innerleftmargin = 1em,
  innerrightmargin = 0pt,
 innerbottommargin = 0pt,
  innertopmargin = 0pt,
  rightmargin = 0pt,
  linewidth = 2pt,
  topline = false,
  rightline = false,
  bottomline = false,
  skipabove = 6pt
  ]{leftbar}


\newcommand{\mframe}[1]{\begin{leftbar}{#1}\end{leftbar}}


%You can copy those commands to the preamble of your document and fill in the values that you prefer (e.g., 0pt for the indents and \baselineskip for the \parskip).








% \titleclass{\subsubsection}{straight}
% \titleformat{\subsubsection}%
%   [hang]% shape
%   {\normalfont\large\itshape}% format applied to label+text
%   {\thesubsubsection}% label
%   {1em}% horizontal separation between label and title body
%   {}% before the title body
%   []% after the title body
  
  
  

%%%%%%%%%%%%%%%%%%%%%%%%%%%%%%%%%%%%%
%%%%%% FANCY FRAMES
%% https://tex.stackexchange.com/questions/348501/example-of-box-inside-box
%%%%%%%%%%%%%%%%%%%%%%%%
%\usepackage[margin=0.5in]{geometry}
%\usepackage{tikz,lipsum,lmodern}
\usepackage{tikz,lipsum}
\usepackage[most]{tcolorbox}

\tcbset{titre/.style={boxed title style={boxrule=0pt,colframe=white}}}

\definecolor{gradientGreenL}{HTML}{1fe2ad} 
\definecolor{gradientGreenR}{HTML}{d4eb6f} 


\newtcolorbox{BoxResume}[2][]{
                boxrule=0.5pt,
                colback=white,
                top=3pt,bottom=2pt,left=2pt,right=2pt,
                colframe=webbrown,
                fonttitle=\sffamily\small,%\bfseries
                coltitle=black,
                colbacktitle=white,
                enhanced,
                attach boxed title to top left={xshift=5mm, yshift=-2mm},
                title=#2,#1
                }


\newtcolorbox{BoxIn}{
enhanced,
colframe=white,
interior style={
left color=gradientGreenL!7!white,
right color=gradientGreenR!7!white},
%frame style image=background\aa.jpg
left=5mm,
top=4mm,
bottom=4mm,
right=5mm,
boxsep=0mm,
nobeforeafter}



\newtcolorbox{BoxResumeNew}[2][]{
                boxrule=1pt,
                colback=white,
                top=3pt,bottom=2pt,left=2pt,right=2pt,
                colframe=black,
                fonttitle=\sffamily\small,%\bfseries
                coltitle=black,
                colbacktitle=white,
                enhanced,
                attach boxed title to top left={xshift=5mm, yshift=-2mm},
                title=#2,#1
                }


\newtcolorbox{BoxInNew}{
enhanced,
colframe=white,
colback=black!2!white,
%frame style image=background\aa.jpg
left=5mm,
top=4mm,
bottom=4mm,
right=5mm,
boxsep=0mm,
nobeforeafter}



\newcommand{\remember}[1]{
\vspace*{\fill}
\begin{BoxResumeNew}[titre]{WHAT YOU MUST REMEMBER}
 \begin{BoxInNew}{}
 #1
 \end{BoxInNew}{}
\end{BoxResumeNew}
\vspace{0.5cm}
} 

%%%% USAGE

%\remember{
%\textit{Contributions:}\vspace{0.5em}
%\begin{itemize}
%\item[$-$] Design and development of a finger robotic actuator for mobile devices
%\item[$-$] Applications and scenarios that demonstrate its use as a medium,  as a tool and as a virtual partner
%\item[$-$] Initial evaluation of perception of the appearance and the relevance of scenarios
%\end{itemize}
%}


%%%%%%%%%%%%%%%%%%%%%%%%%%%%%%
%%%%%%%%%   QUOTE %%%%%%%%%%%%%%%%%%%
%%%%%%%%%%%%%%%%%%

\makeatletter
\renewcommand{\@chapapp}{}% Not necessary...
\newenvironment{chapquote}[2][2em]
  {\setlength{\@tempdima}{#1}%
   \def\chapquote@author{#2}%
   \parshape 1 \@tempdima \dimexpr\textwidth-2\@tempdima\relax%
   }
  {\par\normalfont\hfill--\ \chapquote@author\hspace*{\@tempdima}\par\bigskip}
\makeatother


%\listfiles

% Set the language
\usepackage[english]{babel} % Load characters and hyphenation
\usepackage[english=british]{csquotes} % English quotes

% Load packages for testing
\usepackage{blindtext}
%\usepackage{showframe} % Uncomment to show boxes around the text area, margin, header and footer
%\usepackage{showlabels} % Uncomment to output the content of \label commands to the document where they are used



% Load the bibliography package
\usepackage{styles/kaobiblio}
\addbibresource{main.bib} % Bibliography file

% Load mathematical packages for theorems and related environments. NOTE: choose only one between 'mdftheorems' and 'plaintheorems'.
\usepackage{styles/mdftheorems}
%\usepackage{styles/plaintheorems}

\graphicspath{{examples/documentation/images/}{images/}} % Paths in which to look for images

\makeindex[columns=3, title=Alphabetical Index, intoc] % Make LaTeX produce the files required to compile the index

\makeglossaries % Make LaTeX produce the files required to compile the glossary

\makenomenclature % Make LaTeX produce the files required to compile the nomenclature

% Reset sidenote counter at chapters
%\counterwithin*{sidenote}{chapter}

%----------------------------------------------------------------------------------------

\newcommand{\red}[1]{\textcolor[rgb]{1, 0, 0}{#1}}


\begin{document}

%----------------------------------------------------------------------------------------
%	BOOK INFORMATION
%----------------------------------------------------------------------------------------

% \titlehead{The \texttt{kaobook} class}
% \subject{Course title}
% \title{Project title}
% \subtitle{A very short and explicit subtitle for your project}

% \author{Firstname LASTNAME}

\date{\today}

\titlehead{\centering\includegraphics[width=6cm]{dvic.png}}

% \publishers{\textbf{Supervisior}\\ Clément Duhart or Marc Teyssier or Xiao Xiao }


%----------------------------------------------------------------------------------------

\frontmatter % Denotes the start of the pre-document content, uses roman numerals

%----------------------------------------------------------------------------------------
%	OPENING PAGE
%----------------------------------------------------------------------------------------

% \makeatletter
% \extratitle{
% 	% In the title page, the title is vspaced by 9.5\baselineskip
% 	\vspace*{9\baselineskip}
% 	\vspace*{\parskip}
% 	\begin{center}
% 		% In the title page, \huge is set after the komafont for title
% 		\usekomafont{title}\huge\@title
% 	\end{center}
% }
% \makeatother

%----------------------------------------------------------------------------------------
%	COPYRIGHT PAGE
%----------------------------------------------------------------------------------------

% \makeatletter
% \uppertitleback{\@titlehead} % Header

% \lowertitleback{
% 	\textbf{Disclaimer}\\
% 	You can edit this page to suit your needs. For instance, here we have a no copyright statement, a colophon and some other information. This page is based on the corresponding page of Ken Arroyo Ohori's thesis, with minimal changes.
	
% 	\medskip
	
% 	\textbf{No copyright}\\
% 	\cczero\ This book is released into the public domain using the CC0 code. To the extent possible under law, I waive all copyright and related or neighbouring rights to this work.
	
% 	To view a copy of the CC0 code, visit: \\\url{http://creativecommons.org/publicdomain/zero/1.0/}
	
% 	\medskip
	
% 	\textbf{Colophon} \\
% 	This document was typeset with the help of \href{https://sourceforge.net/projects/koma-script/}{\KOMAScript} and \href{https://www.latex-project.org/}{\LaTeX} using the \href{https://github.com/fmarotta/kaobook/}{kaobook} class.
	
% 	The source code of this book is available at:\\\url{https://github.com/fmarotta/kaobook}
	
% 	(You are welcome to contribute!)
	
% 	\medskip
	
% 	\textbf{Publisher} \\
% 	First printed in May 2019 by \@publishers
% }
% \makeatother

%----------------------------------------------------------------------------------------
%	DEDICATION
%----------------------------------------------------------------------------------------

% \dedication{
% 	The harmony of the world is made manifest in Form and Number, and the heart and soul and all the poetry of Natural Philosophy are embodied in the concept of mathematical beauty.\\
% 	\flushright -- D'Arcy Wentworth Thompson
% }

%----------------------------------------------------------------------------------------
%	OUTPUT TITLE PAGE AND PREVIOUS
%----------------------------------------------------------------------------------------

% Note that \maketitle outputs the pages before here

% If twoside=false, \uppertitleback and \lowertitleback are not printed
% To overcome this issue, we set twoside=semi just before printing the title pages, and set it back to false just after the title pages
\KOMAoptions{twoside=semi}
\maketitle
\KOMAoptions{twoside=false}


%----------------------------------------------------------------------------------------
%	PREFACE
%----------------------------------------------------------------------------------------

\chapter*{\centering \normalsize Preface}
%\addcontentsline{toc}{chapter}{Preface} % Add the preface to the table of contents as a chapter

The De Vinci Innovation Center (DVIC) is a community of makers that develops
technologies within philosophical and critical frameworks to shape our societies’
futures. The objective is to implement real-world solutions as well as design projects
to enhance public engagement, improve education, and overall provide scientific
knowledge. Our researchers contribute actively to top-level international research
in multiple fields, including artificial intelligence, human-computer interactions,
education, and ecology. We believe that these objectives require a transdisciplinary
approach, that bridges the gap between sciences, techniques, sociology, and philosophy. This is performed by collaborating with other scientists and industrial and
startup sharing our values, to form strong research partnerships...



The Artificial Lives group, led by Dr. Clement Duhart, aims to develop the next generation of machines and Human-Machine Interfaces. The group members strongly
believe that through the combination of Design and Engineering, human-centered
technologies can blend into our environments to become invisible, vastly improving
daily lives. To achieve this vision, the members contribute to human-computer
interactions, cognitive enhancement through new forms of extended intelligence,
learning platforms, and cobotic. Our bio-inspired, multidisciplinary approach couples AI and virtual reality with intelligent materials, robotics and the Internet of
Things.



For the past two years, De Vinci Innovation Center (DVIC) students following
the Creative Technologies curriculum had the opportunity to develop their vision
on technology, innovation, and society. This proceeding is a composition of six
master’s theses, ranging from Machine Learning, Human-Computer-Interaction to
Robotics. The authors strongly believe that developing alternative futures requires
new types of engineering that take into consideration both the people’s needs and
the environment. These documents have been written to reflect this vision and
refined over several months with an iterative reviewing supervised by the Principal
Investigators.



The Authors, the Principal Investigators and the whole DVIC community is proud
of releasing this first proceeding. We dedicate this first edition to Pascal Brouaye
and Nelly Rouyres, without whom nothing would have been possible.
 
RF ADD RF PICTURE !!!!!!!!!!!!!!!!!


\pagebreak

\afterpage{\blankpage}
%\index{preface}

%----------------------------------------------------------------------------------------
%	LIST OF THESE
%----------------------------------------------------------------------------------------

\newline
\chapter*{\centering \Huge List of Theses}



\pagebreak

\afterpage{\blankpage}

%----------------------------------------------------------------------------------------
%	TITLE 
%----------------------------------------------------------------------------------------

\vspace{3cm}
\chapter*{\centering \Huge From Human Handwriting to
Robotic Arm Dexterity}
{\centering \textbf Tom Leizerovici}

Motivation: Before dive into this master thesis. For a better comprehension I want to explain a bit my motivation, my vision, my goal about the work I have done during my master. The goal was first learning a lot of new hard skill in a way that it can be used for to create a new vision of what is technology. Particularly In a word where we start to feel the consequence -- and paradoxically the need -- of less energy and resource. 



\pagebreak

%----------------------------------------------------------------------------------------
%	TABLE OF CONTENTS & LIST OF FIGURES/TABLES
%----------------------------------------------------------------------------------------

\begingroup % Local scope for the following commands

% Define the style for the TOC, LOF, and LOT
%\setstretch{1} % Uncomment to modify line spacing in the ToC
% \hypersetup{linkcolor=blue} % Uncomment to set the colour of links in the ToC
\setlength{\textheight}{23cm} % Manually adjust the height of the ToC pages

% Turn on compatibility mode for the etoc package
\etocstandarddisplaystyle % "toc display" as if etoc was not loaded
\etocstandardlines % toc lines as if etoc was not loaded

% Comment both of the following lines to have the LOF and the LOT on different pages
\let\cleardoublepage\bigskip
\let\clearpage\bigskip

%\listoftables % Output the list of tables

\tableofcontents % Output the table of contents






\endgroup

%----------------------------------------------------------------------------------------
%	MAIN BODY
%----------------------------------------------------------------------------------------

\mainmatter % Denotes the start of the main document content, resets page numbering and uses arabic numbers
\setchapterstyle{kao} % Choose the default chapter heading style

%
\section*{Introduction \red{(Phase 1)}}
\textbf{$\sim$1 pages}

The first paragraph describes briefly the project content. 

The second paragraph describes how this project is related with the student's creative technology curriculum strategy (integrated within an innovation track's project or professional perspectives).

Finally, you present here the context and an exhaustive list of related works or existing solution in a narration. This very important task will lead you in making right choices and presenting how your work is important or difficult. 


\pagebreak



%
\section*{Project title \red{(Phase 2)}}
\textbf{$\sim$4 page}

This paragraph must describe briefly the general principle of the project, what it is, how it works and what outputs are expected. Comments if the results are different. 

\subsection*{Overview}
In this section, student must present the general project overview of the system with a brief description of the main components. 


\subsection*{Components}

The student present how are implemented the different components. This section must be very clear with schematics, calculations, models, list of equipment or resources, etc. 

The student creates a new \textit{subsection} for each component. Although the grading is slightly different between the Designer and Engineer, we encourage the students to give as many details and many subsections as possible (both from Engineer and Design). It is also possible to add other subsections. 

\textbf{Project components for Creative Technology Engineer}

Here is a list of subsection that and Engineer CreaTech must present (at least 4). 

\begin{itemize}
    \item System Architecture
    \item Simulation and / or Modelization
    \item Fabrication and / or Implementation
\end{itemize}


\textbf{Project components for Creative Technology Designer}

Here is a list of subsection that and Design. CreaTech must describe (at least 4). 


\begin{itemize}
    \item 3D mockup and / or Design 
    \item Ergonomy and / or Finish Quality
    \item Fabrication
\end{itemize}

\subsection*{Product Specific Success Criteria}

The student proposes a list of success criteria specific of the projects with the corresponding evaluation marks.

\textbf{Example of PSSC}
\begin{table}[h!]
    \begin{tabular}{|c|l|c|}
      \hline
      Items & Description & Mark (total 20) \\
      \hline
      Robot Navigation & The robot moves forward a straight line & 1 \\
      \cline{2-3} & The robot avoids obstacles  & 3 \\
      \cline{2-3} & The robot plans trajectory obstacles  & 2 \\
      \hline 
      Robot Arm & The robot arm is moving & 1 \\
      \hline 
      ... & ... & ... \\
      \hline 
      Total &  & 20 \\
      \hline
    \end{tabular}
\end{table}


\pagebreak

%

\section*{Evaluation and User Studies \red{(Phase 3)}}
\textbf{$\sim$3 page}







\textbf{Evaluation expectation for CreaTech Engineer (ESILV)}

 This section includes one or several \textit{quantitative} evaluation (\textit{hardware experimentation}). 
 
 Describe the goal of the evaluation, your protocol, and the results with graphs. Comments the results. How different is your work from related work or other similar projects ? What did you prove with this evaluation ? 
 
\textbf{Evaluation expectation for CreaTech Designer (IIM)}

 This section includes a  \textit{qualitative} evaluation, or  \textit{user study}. 
 
 Describe the goal of the user study, what you want to learn from the users interacting or using your system. Evaluate through qualitative questionnaire the user experience of the system or an interaction technique. Describe how you defined the questions and the scales you used. 

Draw inspiration from the evaluation section from the HCI papers. Report the results using graphs or likert scales items. 



%
\section*{Conclusion \red{(Phase 3)}}
\textbf{$\sim$1 page }


This section should briefly conclude the document. It should sum-up the project, the context, and the evaluations.

\pagebreak

\setchapterpreamble[u]{\margintoc}
\chapter{Introduction}
\labch{intro}

\section{intro}
\section{intro2}
\section{intro3}
\subssection{into3.1?}


\pagebreak



%----------------------------------------------------------------------------------------

\backmatter % Denotes the end of the main document content
\setchapterstyle{plain} % Output plain chapters from this point onwards

%----------------------------------------------------------------------------------------
%	BIBLIOGRAPHY
%----------------------------------------------------------------------------------------

% The bibliography needs to be compiled with biber using your LaTeX editor, or on the command line with 'biber main' from the template directory

\defbibnote{bibnote}{Here are the references in citation order.\par\bigskip} % Prepend this text to the bibliography
\printbibliography[heading=bibintoc, title=References, prenote=bibnote] % Add the bibliography heading to the ToC, set the title of the bibliography and output the bibliography note

%----------------------------------------------------------------------------------------
%	NOMENCLATURE
%----------------------------------------------------------------------------------------

% The nomenclature needs to be compiled on the command line with 'makeindex main.nlo -s nomencl.ist -o main.nls' from the template directory

% \nomenclature{$c$}{Speed of light in a vacuum inertial frame}
% \nomenclature{$h$}{Planck constant}

% \renewcommand{\nomname}{Notation} % Rename the default 'Nomenclature'
% \renewcommand{\nompreamble}{The next list describes several symbols that will be later used within the body of the document.} % Prepend this text to the nomenclature

% \printnomenclature % Output the nomenclature

%----------------------------------------------------------------------------------------
%	GREEK ALPHABET
% 	Originally from https://gitlab.com/jim.hefferon/linear-algebra
%----------------------------------------------------------------------------------------

% \vspace{1cm}

% {\usekomafont{chapter}Greek Letters with Pronounciation} \\[2ex]
% \begin{center}
% 	\newcommand{\pronounced}[1]{\hspace*{.2em}\small\textit{#1}}
% 	\begin{tabular}{l l @{\hspace*{3em}} l l}
% 		\toprule
% 		Character & Name & Character & Name \\ 
% 		\midrule
% 		$\alpha$ & alpha \pronounced{AL-fuh} & $\nu$ & nu \pronounced{NEW} \\
% 		$\beta$ & beta \pronounced{BAY-tuh} & $\xi$, $\Xi$ & xi \pronounced{KSIGH} \\ 
% 		$\gamma$, $\Gamma$ & gamma \pronounced{GAM-muh} & o & omicron \pronounced{OM-uh-CRON} \\
% 		$\delta$, $\Delta$ & delta \pronounced{DEL-tuh} & $\pi$, $\Pi$ & pi \pronounced{PIE} \\
% 		$\epsilon$ & epsilon \pronounced{EP-suh-lon} & $\rho$ & rho \pronounced{ROW} \\
% 		$\zeta$ & zeta \pronounced{ZAY-tuh} & $\sigma$, $\Sigma$ & sigma \pronounced{SIG-muh} \\
% 		$\eta$ & eta \pronounced{AY-tuh} & $\tau$ & tau \pronounced{TOW (as in cow)} \\
% 		$\theta$, $\Theta$ & theta \pronounced{THAY-tuh} & $\upsilon$, $\Upsilon$ & upsilon \pronounced{OOP-suh-LON} \\
% 		$\iota$ & iota \pronounced{eye-OH-tuh} & $\phi$, $\Phi$ & phi \pronounced{FEE, or FI (as in hi)} \\
% 		$\kappa$ & kappa \pronounced{KAP-uh} & $\chi$ & chi \pronounced{KI (as in hi)} \\
% 		$\lambda$, $\Lambda$ & lambda \pronounced{LAM-duh} & $\psi$, $\Psi$ & psi \pronounced{SIGH, or PSIGH} \\
% 		$\mu$ & mu \pronounced{MEW} & $\omega$, $\Omega$ & omega \pronounced{oh-MAY-guh} \\
% 		\bottomrule
% 	\end{tabular} \\[1.5ex]
% 	Capitals shown are the ones that differ from Roman capitals.
% \end{center}

%----------------------------------------------------------------------------------------
%	GLOSSARY
%----------------------------------------------------------------------------------------

% The glossary needs to be compiled on the command line with 'makeglossaries main' from the template directory

% \newglossaryentry{computer}{
% 	name=computer,
% 	description={is a programmable machine that receives input, stores and manipulates data, and provides output in a useful format}
% }

% Glossary entries (used in text with e.g. \acrfull{fpsLabel} or \acrshort{fpsLabel})
\newacronym[longplural={Frames per Second}]{fpsLabel}{FPS}{Frame per Second}
\newacronym[longplural={Tables of Contents}]{tocLabel}{TOC}{Table of Contents}

\setglossarystyle{listgroup} % Set the style of the glossary (see https://en.wikibooks.org/wiki/LaTeX/Glossary for a reference)
\printglossary[title=Special Terms, toctitle=List of Terms] % Output the glossary, 'title' is the chapter heading for the glossary, toctitle is the table of contents heading

%----------------------------------------------------------------------------------------
%	INDEX
%----------------------------------------------------------------------------------------

% The index needs to be compiled on the command line with 'makeindex main' from the template directory

\printindex % Output the index

%----------------------------------------------------------------------------------------
%	BACK COVER
%----------------------------------------------------------------------------------------

% If you have a PDF/image file that you want to use as a back cover, uncomment the following lines

%\clearpage
%\thispagestyle{empty}
%\null%
%\clearpage
%\includepdf{cover-back.pdf}

%----------------------------------------------------------------------------------------

\usepackage{afterpage}

\newcommand\blankpage{%
    \null
    \thispagestyle{empty}%
    %\addtocounter{page}{-1}%
    \newpage}



\end{document}
